%======================================================================
\section{Related work}
\label{sec:related}

The ingredients of this paper are classical, and it is worth saying
precisely which shelf each came from and what, if anything, is new in
the assembly.

\emph{Densities from option prices.}  That the call curve's second
strike derivative is the risk-neutral density is Breeden and
Litzenberger's \citep{BreedenLitzenberger1978}; the static no-arbitrage
characterization of call quotes used in \cref{def:slice} is classical
and is stated compactly by Carr and Madan \citep{CarrMadan2005}.  This
paper's \eqref{eq:bl_division} is the same identity read in the model's
coordinates, where it costs one division.

\emph{Wing asymptotics.}  The moment formula bounding wing slopes is
Lee's \citep{Lee2004}; the upgrade from $\limsup$ to limits under
regular variation is Benaim and Friz's \citep{BenaimFriz2009}.  The
contribution here is only the closed composition
\eqref{eq:leeclosed} --- the model's tail coordinates happen to make
Lee's formula an explicit two-parameter map --- and the effective-slope
honesty figure.

\emph{Implied-volatility parametrizations.}  SVI
\citep{Gatheral2004} and its surface extension SSVI
\citep{GatheralJacquier2014} are the industry benchmarks for the
volatility chart, with explicitly characterized (but parameter-coupled)
no-arbitrage regions; Fengler's arbitrage-free smoothing
\citep{Fengler2009} represents the spline-with-constraints tradition in
the price chart.  Both approaches place validity in the optimizer;
the present model places it in the coordinates.  The trade is not free
--- SVI's five parameters are quotable desk vocabulary in a way sixteen
Legendre coefficients are not --- and \cref{sec:atm} plus the package
machinery of \cref{app:implementation} are the model's answer, not a
dismissal of the need.

\emph{Quantile-based modeling.}  Modeling distributions through
quantile functions goes back at least to Parzen \citep{Parzen1979} and
is developed monographically by Gilchrist \citep{Gilchrist2000}; Keelin's
metalog family \citep{Keelin2016} builds flexible unbounded
distributions as polynomial expansions of the quantile function itself
and is the closest relative in the applied literature.  The two design
differences are load-bearing: LQD expands the \emph{logarithm} of the
quantile density --- so monotonicity is structural at every order,
where the metalog must verify feasibility of its coefficient vector ---
and it carries the $-\log u-\log(1-u)$ skeleton explicitly, so the
expansion coefficients control tails through two endpoint values with
the martingale condition solvable by one shift.  In functional data
analysis, Petersen and M\"uller \citep{PetersenMuller2016} use the same
log-quantile-density transform to map density-valued data to an
unconstrained Hilbert space --- the identical unconstrainedness
argument, made for statistical rather than arbitrage reasons; this
paper adds the financial boundary conditions (martingale normalization,
the $\lambda_+<1$ wall, Lee consistency) and the pricing calculus.

\emph{Variance swaps and the log contract.}  The log-contract
replication is Neuberger's \citep{Neuberger1994}, its dealer-standard
exposition Demeterfi, Derman, Kamal, and Zou's
\citep{DemeterfiDermanKamalZou1999}, and the general payoff-replication
frame Carr and Madan's \citep{CarrMadan1998}.  Here the fitted law is
explicit, so \eqref{eq:varswap} integrates the density directly rather
than replicating across a strip.

\emph{Calendar order.}  That call-price domination with equal means is
convex order, and that convex-ordered marginals are exactly those a
martingale can traverse, is Strassen's and Kellerer's
\citep{Strassen1965,Kellerer1972}.  \Cref{thm:ledgerorder} restates the
first fact as a pointwise comparison of the model's native
computational object; the conjugate-duality proof is elementary and
surely folklore in some notation, but the author has not found the
ledger form stated.

\emph{Numerics.}  Monotone cubic interpolation certificates are Fritsch
and Carlson's \citep{FritschCarlson1980}; floating-point discipline
follows Higham \citep{Higham2002}; the solver is the trust-region
reflective method \citep{BranchColemanLi1999} as implemented in SciPy
\citep{SciPy2020}.

What, then, is claimed as new?  Not the transform, not the moment
formula, not the conjugacy of calls and tail expectations.  The claims
are architectural: the separation of validity, information, and
convention into structure, audit, and disclosed choices respectively;
the single-ledger organization in which pricing, differentiation
(\cref{thm:envelope}), density, variance swap, and calendar order are
five reads of one array; the closed ATM microscope \eqref{eq:atm_w}
with the skew-as-digital-mismatch reading; and the certificate-plus-audit
doctrine of \cref{sec:computation} by which the discretization earns
the continuous theorems.  Readers are invited to judge the assembly,
since every part has prior art somewhere.
